% =====================================================================
% CHAPTER 7: PIPELINE III — UNSUPERVISED FEATURE EXTRACTION VIA
%             VARIATIONAL AND DENOISING AUTOENCODERS
% =====================================================================
\chapter{Pipeline III: Unsupervised Feature Extraction via Variational and Denoising Autoencoders}
\label{ch:autoencoder}

\section{Introduction and Research Questions}

The preceding chapters have explored fraud detection through graph-based relational reasoning (Chapter~\ref{ch:gnn}) and biologically inspired spiking computation (Chapter~\ref{ch:snn}), both of which operate within the supervised or semi-supervised learning paradigm where labeled fraud examples are available during training. However, the fundamental premise of supervised fraud detection harbours a critical vulnerability: the inductive bias of gradient-optimized classifiers is shaped overwhelmingly by the majority class when the class distribution is as extreme as 98.90\% legitimate versus 1.10\% fraudulent. Under such conditions, decision boundaries are effectively drawn to separate the dense cloud of legitimate transactions from an almost negligible number of scattered fraud instances, and the resulting classifiers exhibit brittle performance---small perturbations in feature space can cause dramatic swings in fraud detection rates, while the sheer volume of legitimate transactions ensures that even a modest false positive rate translates into an unmanageable number of false alerts in production.

This inductive bias failure motivates a fundamentally different approach: rather than learning to discriminate fraud from legitimacy directly, we ask whether the structure of legitimate transactions alone can be captured so precisely that deviations from this structure serve as reliable fraud indicators. Autoencoders---neural networks trained to reconstruct their inputs through a compressed latent bottleneck---provide a natural framework for this approach. When trained exclusively on legitimate transactions, the autoencoder learns a manifold that approximates the data distribution of normal behaviour. Fraudulent transactions, which by definition deviate from this normal manifold, should exhibit elevated reconstruction errors, and these errors can be leveraged as unsupervised anomaly scores without ever exposing the model to a single fraudulent example during training.

Yet the application of autoencoders to tabular fraud detection is far from straightforward. Vanilla autoencoders suffer from the identity mapping problem, where the network learns an approximately invertible transformation that reconstructs all inputs---including anomalies---with low error, thereby collapsing the anomaly signal. Variational Autoencoders (VAEs)~\cite{kingma2014auto} address this issue by imposing a probabilistic structure on the latent space through a Gaussian prior, which regularizes the encoder distribution and theoretically prevents the network from memorizing arbitrary input-output mappings. Denoising Autoencoders (DAEs)~\cite{vincent2010stacked} adopt a complementary strategy, corrupting the input with stochastic noise and training the network to recover the original signal, thereby forcing the model to learn robust features that capture the underlying data manifold rather than superficial pixel-level correspondences. The question of whether these regularisation mechanisms translate into effective anomaly detection under the extreme class imbalance characteristic of financial fraud datasets remains open.

This chapter addresses these gaps through the following research questions:

\begin{description}
\item[RQ1:] Can the probabilistic latent space imposed by a Variational Autoencoder mitigate posterior collapse and improve anomaly detection performance on highly imbalanced tabular fraud data, compared to deterministic reconstruction-based alternatives?
\item[RQ2:] Does stochastic feature corruption via tabular swap noise preserve marginal feature statistics while forcing the Denoising Autoencoder to learn invariant representations that amplify anomalous reconstruction errors for out-of-distribution fraud instances?
\item[RQ3:] Does score-level stacking---combining unsupervised autoencoder anomaly features with gradient-boosted tree meta-learners---yield performance improvements over both standalone autoencoder scoring and purely supervised baselines, and if so, what is the magnitude of this hybridization benefit?
\end{description}

To investigate these questions, we design and implement a complete autoencoder-based fraud detection pipeline featuring both a Variational Autoencoder with Gaussian prior and reparameterization trick, and a Denoising Autoencoder with tabular swap noise at a 15\% corruption rate. The pipeline incorporates a target manifold isolation strategy---training exclusively on legitimate transactions---and a score-level stacking mechanism that concatenates latent coordinates, element-wise reconstruction errors, and original features as input to XGBoost~\cite{chen2016xgboost} and LightGBM~\cite{ke2017lightgbm} meta-learners. The entire pipeline is evaluated on the Bank Account Fraud (BAF) Dataset~\cite{feedzai2022baf,jesus2022bank}, NeurIPS 2022 Variant~I, comprising approximately one million instances with a fraud prevalence of 1.10\%. The results reveal a striking dichotomy: standalone autoencoder anomaly scores achieve near-random discrimination (AUC~$\approx$~0.52--0.55), yet stacked hybrid models reach AUC values of 0.87--0.89, demonstrating that unsupervised features, while insufficient in isolation, provide complementary information that substantially enhances supervised classifiers.

\section{Theoretical Background}

\subsection{Mathematical Foundations of Autoassociative Networks}

Autoencoders belong to the family of autoassociative neural networks, which learn to map an input vector to itself through a constrained intermediate representation. Formally, an autoencoder consists of two parametric functions: an encoder $g_\phi \colon \mathbb{R}^D \to \mathbb{R}^d$ that maps the input $\mathbf{x} \in \mathbb{R}^D$ to a latent code $\mathbf{z} = g_\phi(\mathbf{x}) \in \mathbb{R}^d$, and a decoder $f_\theta \colon \mathbb{R}^d \to \mathbb{R}^D$ that reconstructs the input from the latent representation $\hat{\mathbf{x}} = f_\theta(\mathbf{z})$. The bottleneck constraint $d \ll D$ forces the network to compress the input into a lower-dimensional representation that captures the most salient features of the data distribution. The parameters $\phi$ and $\theta$ are jointly optimized to minimize the reconstruction error, typically measured using the mean squared error (MSE) loss:

\begin{equation}
\mathcal{L}_{\text{AE}}(\phi, \theta) = \frac{1}{N} \sum_{i=1}^{N} \bigl\| \mathbf{x}_i - f_\theta\bigl(g_\phi(\mathbf{x}_i)\bigr) \bigr\|^2 = \frac{1}{N} \sum_{i=1}^{N} \bigl\| \mathbf{x}_i - \hat{\mathbf{x}}_i \bigr\|^2
\label{eq:ae-mse-loss}
\end{equation}

The theoretical appeal of autoencoders for anomaly detection rests on the manifold hypothesis: if the training data lies on or near a low-dimensional manifold $\mathcal{M} \subset \mathbb{R}^D$, then the autoencoder should learn to reconstruct inputs that lie on this manifold with low error, while inputs that lie far from the manifold---anomalies---should exhibit high reconstruction error. This discrepancy forms the basis for anomaly scoring: samples with large $\|\mathbf{x} - \hat{\mathbf{x}}\|^2$ are flagged as potential anomalies.

However, the vanilla autoencoder suffers from a fundamental limitation known as the identity mapping problem. When the encoder and decoder are sufficiently expressive---which is the case for modern deep architectures with multiple hidden layers and nonlinear activations---the network can learn an approximately invertible transformation that maps any input to a latent code and back with minimal reconstruction error, regardless of whether the input lies on the data manifold. This is because the composition $f_\theta \circ g_\phi$ can approximate the identity function on the entire input space $\mathbb{R}^D$, not just on the data manifold $\mathcal{M}$. In the context of fraud detection, this means that fraudulent transactions can be reconstructed nearly as accurately as legitimate ones, collapsing the anomaly signal and rendering the reconstruction error an ineffective discriminator. The severity of this problem is compounded in tabular settings where features are not spatially correlated (unlike images), making it easier for the network to learn independent per-feature reconstructions that bypass the manifold-learning objective.

Several strategies have been proposed to mitigate the identity mapping problem, including bottleneck constraints, sparse regularisation, contractive penalties, and the two approaches central to this chapter: probabilistic latent regularisation (VAE) and stochastic input corruption (DAE). Each of these strategies introduces an inductive bias that prevents the network from learning a trivial identity mapping and instead forces it to capture the statistical structure of the training data in a way that generalises to unseen samples from the same distribution but fails for out-of-distribution inputs.

\subsection{Denoising Autoencoders and Marginal-Preserving Tabular Corruption}

The Denoising Autoencoder (DAE), introduced by Vincent et al.~\cite{vincent2010stacked}, addresses the identity mapping problem by modifying the reconstruction objective: rather than learning to map $\mathbf{x}$ to $\mathbf{x}$, the DAE learns to map a corrupted version $\tilde{\mathbf{x}}$ back to the original clean input $\mathbf{x}$. This seemingly simple modification has profound implications for the learned representation. By forcing the network to recover the original signal from a degraded version, the DAE must learn features that capture the statistical dependencies between input dimensions---only by understanding the joint distribution of the features can the network infer the values of corrupted dimensions from the uncorrupted ones. The DAE objective is:

\begin{equation}
\mathcal{L}_{\text{DAE}}(\phi, \theta) = \frac{1}{N} \sum_{i=1}^{N} \bigl\| \mathbf{x}_i - f_\theta\bigl(g_\phi(\tilde{\mathbf{x}}_i)\bigr) \bigr\|^2
\label{eq:dae-loss}
\end{equation}

where $\tilde{\mathbf{x}}_i$ is the corrupted version of $\mathbf{x}_i$ obtained through a stochastic corruption process $q_D(\tilde{\mathbf{x}} \mid \mathbf{x})$.

For image data, the most common corruption strategy is additive Gaussian noise or salt-and-pepper noise, which introduces random perturbations to pixel values. However, these corruption mechanisms are poorly suited for tabular data, where features are heterogeneous (binary, categorical, continuous) and the notion of spatial smoothness that justifies Gaussian noise in images does not apply. Instead, we employ tabular swap noise, a corruption mechanism specifically designed for structured tabular data. For each feature $j$ in a sample $\mathbf{x}_i$, swap noise replaces the feature value with a value drawn from the marginal distribution of that feature across the training set:

\begin{equation}
\tilde{x}_{i,j} = \begin{cases}
x_{i,j} & \text{with probability } 1 - \nu \\
x_{k,j} \sim \text{Marginal}(x_j) & \text{with probability } \nu
\end{cases}
\label{eq:dae-swap-noise}
\end{equation}

where $\nu = 0.15$ is the corruption rate and $x_{k,j}$ is the value of feature $j$ from a randomly selected training sample $k$. The swap noise mechanism has several desirable properties for tabular data. First, it preserves the marginal distribution of each feature: the corrupted value $\tilde{x}_{i,j}$ is always a legitimate value observed in the training data, ensuring that the DAE never encounters out-of-range inputs during training. Second, it breaks the joint distribution of features: when feature $j$ is swapped, its value is drawn independently from other features, destroying the dependency structure that the DAE must then learn to reconstruct. Third, the corruption rate $\nu$ controls the difficulty of the denoising task: higher corruption rates force the network to rely more heavily on cross-feature dependencies to reconstruct the original input, leading to richer feature representations but potentially higher variance in the reconstruction.

The architecture adopted for the DAE follows a symmetric hourglass design: $128 \to 64 \to 32 \to 64 \to 128 \to 53$. The encoder compresses the 53-dimensional input through progressively narrower hidden layers (128, 64, 32), creating a 32-dimensional bottleneck representation. The decoder then expands this representation back through layers of increasing width (64, 128), culminating in a 53-dimensional output that reconstructs the original input. The output dimension of 53 corresponds to the expanded feature space after one-hot encoding of categorical variables. ReLU activations are applied to all hidden layers, while the output layer uses a sigmoid activation to ensure reconstructed values lie in the normalized range $[0, 1]$. The target-isolated loss function ensures that the DAE is trained exclusively on legitimate transactions, establishing the ``normal manifold'' that fraudulent transactions are expected to deviate from.

\subsection{Variational Autoencoders and Probabilistic Manifold Topologies}

The Variational Autoencoder (VAE), introduced by Kingma and Welling~\cite{kingma2014auto}, takes a fundamentally different approach to regularising the latent space. Rather than deterministically mapping each input to a single latent code, the VAE encodes each input as a probability distribution over the latent space. The encoder outputs the parameters of a Gaussian distribution---a mean vector $\boldsymbol{\mu}_\phi(\mathbf{x})$ and a diagonal covariance matrix $\text{diag}(\boldsymbol{\sigma}^2_\phi(\mathbf{x}))$---and the latent code is sampled from this distribution:

\begin{equation}
q_\phi(\mathbf{z} \mid \mathbf{x}) = \mathcal{N}\bigl(\boldsymbol{\mu}_\phi(\mathbf{x}),\; \text{diag}(\boldsymbol{\sigma}^2_\phi(\mathbf{x}))\bigr)
\label{eq:vae-encoder}
\end{equation}

The stochastic sampling from $q_\phi(\mathbf{z} \mid \mathbf{x})$ introduces a fundamental challenge for gradient-based optimization: the sampling operation is non-differentiable, preventing the backpropagation of gradients through the latent code. The reparameterization trick resolves this by reformulating the sampling as a differentiable transformation of a noise variable:

\begin{equation}
\mathbf{z} = \boldsymbol{\mu}_\phi(\mathbf{x}) + \boldsymbol{\sigma}_\phi(\mathbf{x}) \odot \boldsymbol{\epsilon}, \quad \boldsymbol{\epsilon} \sim \mathcal{N}(\mathbf{0}, \mathbf{I})
\label{eq:vae-reparameterization}
\end{equation}

where $\odot$ denotes element-wise multiplication and $\boldsymbol{\epsilon}$ is an auxiliary noise variable sampled from a standard normal distribution. Since $\boldsymbol{\epsilon}$ is sampled independently of the network parameters, gradients can flow through $\boldsymbol{\mu}_\phi$ and $\boldsymbol{\sigma}_\phi$ unimpeded, enabling end-to-end optimization.

The VAE is trained by maximizing the Evidence Lower Bound (ELBO), which decomposes into a reconstruction term and a regularisation term:

\begin{equation}
\mathcal{L}_{\text{ELBO}} = \underbrace{-\mathbb{E}_{q_\phi(\mathbf{z} \mid \mathbf{x})}\bigl[\log p_\theta(\mathbf{x} \mid \mathbf{z})\bigr]}_{\text{Reconstruction loss}} + \underbrace{D_{\text{KL}}\bigl(q_\phi(\mathbf{z} \mid \mathbf{x}) \,\|\, p(\mathbf{z})\bigr)}_{\text{Regularisation term}}
\label{eq:vae-elbo}
\end{equation}

The reconstruction term encourages the decoder to accurately reconstruct the input from the sampled latent code, measured via the negative log-likelihood under the decoder distribution. For continuous inputs, this is typically implemented as the mean squared error or the negative log-likelihood of a Gaussian decoder:

\begin{equation}
\log p_\theta(\mathbf{x} \mid \mathbf{z}) = -\frac{1}{2} \sum_{j=1}^{D} \left(\frac{x_j - \hat{x}_j}{\sigma^2}\right) - \frac{D}{2}\log(2\pi\sigma^2)
\label{eq:vae-nll}
\end{equation}

The KL divergence term penalises deviations of the encoder distribution from the prior $p(\mathbf{z}) = \mathcal{N}(\mathbf{0}, \mathbf{I})$, a standard normal distribution. For the Gaussian encoder with diagonal covariance, this term admits a closed-form solution:

\begin{equation}
D_{\text{KL}}\bigl(q_\phi(\mathbf{z} \mid \mathbf{x}) \,\|\, p(\mathbf{z})\bigr) = -\frac{1}{2} \sum_{j=1}^{d} \left(1 + \log \sigma_{\phi,j}^2 - \mu_{\phi,j}^2 - \sigma_{\phi,j}^2\right)
\label{eq:vae-kl}
\end{equation}

The KL divergence term serves as a probabilistic regulariser that prevents the latent space from developing isolated clusters corresponding to individual training examples---a phenomenon known as posterior collapse, where the encoder distributions become increasingly peaked (small $\sigma$) and the latent codes degenerate into a discrete codebook. By encouraging the encoder distributions to remain close to the isotropic Gaussian prior, the KL term ensures that the latent space is continuous and smooth, with neighbouring latent codes corresponding to semantically similar inputs. This smoothness property is particularly important for anomaly detection: if the latent space is well-regularised, then anomalous inputs should map to regions of the latent space that are far from the high-density regions occupied by legitimate transactions, and the decoder should produce poor reconstructions for these out-of-distribution latent codes.

However, the tension between the reconstruction and regularisation objectives---often called the KL--reconstruction trade-off---can lead to a different form of posterior collapse where the KL term dominates and the encoder learns to ignore the input entirely, mapping all inputs to the prior distribution. This ``over-regularised'' regime produces a decoder that generates average-looking outputs regardless of the input, destroying the anomaly detection signal. Balancing these two objectives is a central challenge in applying VAEs to anomaly detection, and we revisit this issue in the experimental results section.

\subsection{Information Fusion Hybridization and Downstream Stacking Loops}

The core insight driving the stacking architecture in this pipeline is that unsupervised autoencoder features and supervised gradient-boosted features capture fundamentally different aspects of the fraud detection problem. Autoencoders learn to model the distribution of normal transactions; when a fraudulent transaction deviates from this normal manifold, the resulting reconstruction error provides a distributional anomaly signal that is complementary to the discriminative signal captured by supervised classifiers. However, as our experiments demonstrate, this anomaly signal is too weak and noisy to serve as a standalone detector. The stacking framework resolves this by using autoencoder-derived features as additional inputs to a supervised meta-learner, which can learn to weight and combine these features optimally.

The feature fusion strategy proceeds as follows. For each transaction $\mathbf{x}_i$, the trained autoencoder produces two types of derived features: the latent coordinates $\mathbf{z}_i = g_\phi(\mathbf{x}_i) \in \mathbb{R}^d$, which capture the compressed representation of the transaction in the learned manifold, and the element-wise absolute reconstruction errors $\boldsymbol{\epsilon}_i = |\mathbf{x}_i - \hat{\mathbf{x}}_i| \in \mathbb{R}^D$, which quantify the per-feature discrepancy between the original input and its reconstruction. These derived features are concatenated with the original feature vector through a concatenation hub:

\begin{equation}
\mathbf{h}_i = \bigl[\,\mathbf{x}_i \;\|\; \mathbf{z}_i \;\|\; \boldsymbol{\epsilon}_i\,\bigr] \in \mathbb{R}^{D + d + D}
\label{eq:ae-concat-hub}
\end{equation}

where $\|$ denotes concatenation. The resulting augmented feature vector $\mathbf{h}_i$ has dimension $D + d + D = 53 + d + 53$, providing the meta-learner with access to the original features, the compressed latent representation, and the anomaly signal simultaneously. The meta-learner---either XGBoost~\cite{chen2016xgboost} or LightGBM~\cite{ke2017lightgbm}---then learns a discriminative function on this augmented feature space, effectively performing a form of learned feature weighting that determines how much to rely on the unsupervised anomaly signal versus the raw features for each prediction.

The choice of gradient-boosted decision trees as the meta-learner is motivated by several considerations. First, tree-based ensembles are known to be the strongest performers on tabular data~\cite{grinsztajn2022tree}, where they consistently outperform deep learning approaches on heterogeneous feature sets. Second, gradient-boosted trees can naturally handle the mixed-scale nature of the concatenated feature vector, which contains raw features, low-dimensional latent codes, and per-feature error magnitudes, without the feature normalisation that neural networks would require. Third, the sequential nature of boosting allows the model to focus on the most informative features first, automatically down-weighting uninformative autoencoder-derived features while preserving the useful anomaly signals. Fourth, the non-parametric nature of tree-based models avoids the distributional assumptions that could limit the effectiveness of linear meta-learners.

\section{Deep Feature Space Interaction and Regularization Analysis}

\subsection{Decision Tree Partitioning Under Extreme Class Imbalance}

Gradient-boosted decision trees partition the feature space through a sequence of axis-aligned splits, each selected to maximize the reduction in impurity as measured by the Gini index or cross-entropy. For a node containing $N$ samples with class proportions $p_1, p_2, \ldots, p_K$, the Gini impurity is defined as:

\begin{equation}
G = 1 - \sum_{k=1}^{K} p_k^2
\label{eq:ae-gini}
\end{equation}

Under extreme class imbalance ($p_{\text{legit}} \approx 0.989$, $p_{\text{fraud}} \approx 0.011$), the Gini impurity of the root node is $G \approx 2 \times 0.989 \times 0.011 \approx 0.0218$, which is already very close to zero. This near-purity of the root node means that most splits produce only marginal reductions in impurity, and the resulting trees tend to be shallow and uninformative. More critically, the few fraud samples in each node are vulnerable to being treated as noise---a single mislabeled or atypical legitimate transaction can have a disproportionate influence on the split criterion, causing the tree to overfit to spurious patterns in the minority class rather than learning the true decision boundary.

This phenomenon is exacerbated by the feature-space geometry of fraud detection: fraudulent transactions are not cleanly separated from legitimate ones in feature space but instead lie in the tails and intersections of the legitimate distribution. Under these conditions, decision trees tend to carve out small, isolated regions around individual fraud samples, creating high-variance models that generalize poorly to unseen fraud patterns. The autoencoder-derived features address this problem by providing an alternative characterisation of the data manifold: rather than relying solely on the original features to partition the space, the meta-learner can leverage the reconstruction error profile $\boldsymbol{\epsilon}_i$ as a manifold-proximity signal that captures whether a transaction lies near or far from the learned normal manifold. This additional signal can guide the tree splits toward regions where the autoencoder detects distributional anomalies, complementing the discriminative information captured by the original features.

\subsection{Stochastic Forcing Regularizer Mechanics}

The swap noise mechanism in the DAE functions as a multi-feature regularizer whose effects extend beyond the simple prevention of identity mapping. When a feature $j$ is replaced by a value drawn from its marginal distribution, the network must infer the original value of $x_j$ from the remaining uncorrupted features. This inference requires the network to learn the conditional distribution $p(x_j \mid \mathbf{x}_{\setminus j})$, where $\mathbf{x}_{\setminus j}$ denotes all features except $j$. Since the conditional distribution depends on the joint distribution of all features, the DAE is implicitly trained to model the full joint distribution of the training data through its denoising objective.

The corruption rate $\nu = 0.15$ means that, on average, approximately 8 of the 53 features in each sample are corrupted during training. With multiple features simultaneously corrupted, the network must perform a form of multi-variate imputation, recovering the original values of all corrupted features from the partially observed signal. This multi-feature corruption regime forces the network to learn robust, cross-feature dependencies rather than relying on any single feature to predict its own value. The stochastic nature of the corruption---different features are corrupted in different samples---ensures that the network cannot develop feature-specific shortcuts and must instead learn a general-purpose representation of the data manifold.

From a regularisation perspective, swap noise can be understood as a data augmentation strategy that generates an infinite variety of partially corrupted inputs from the finite training set. Each training epoch presents a different set of corrupted inputs, preventing the network from memorizing the training data and forcing it to learn the statistical structure that generalises across corrupted and uncorrupted inputs. This implicit regularisation is particularly valuable in the fraud detection setting, where the training set (legitimate transactions only) is large but the risk of overfitting to the specific characteristics of the training period remains a concern due to temporal distribution shift.

\subsection{Element-Wise Absolute Errors as Feature Attention}

The element-wise absolute reconstruction error $\epsilon_{i,j} = |x_{i,j} - \hat{x}_{i,j}|$ serves a dual role in the stacking framework. At the sample level, the aggregate reconstruction error $\|\boldsymbol{\epsilon}_i\|$ provides a scalar anomaly score that quantifies the overall deviation of a transaction from the learned manifold. At the feature level, the individual error components $\epsilon_{i,j}$ provide a feature-wise anomaly profile that indicates which specific features of a transaction are most poorly reconstructed.

This feature-wise decomposition is particularly informative for fraud detection because fraudulent transactions often exhibit anomalous values on a subset of features rather than globally across all features. For example, a fraudulent application might have a normal-looking credit score and employment status but an anomalous income-to-loan ratio or an unusually recent account opening date. The element-wise error vector $\boldsymbol{\epsilon}_i$ naturally highlights these feature-specific anomalies, providing the meta-learner with a fine-grained attention signal that directs its focus toward the most suspicious features of each transaction.

The use of absolute errors rather than squared errors is a deliberate design choice. While squared errors penalise large deviations more heavily, they also amplify the influence of outlier features that happen to have high reconstruction error for legitimate transactions (e.g., features with high variance or heavy-tailed distributions). Absolute errors provide a more robust measure of reconstruction quality that is less sensitive to these outlier features, ensuring that the anomaly signal is dominated by genuine distributional anomalies rather than reconstruction artefacts. This robustness is crucial in the stacking framework, where the meta-learner must be able to distinguish between reconstruction errors that indicate fraud and those that are merely artefacts of the autoencoder's imperfect reconstruction of high-variance features.

\begin{figure}[H]
\centering
\includegraphics[width=0.92\textwidth]{ae_fig_00}
\caption{Architecture of the autoencoder-based fraud detection pipeline. The left panel shows the VAE pathway with probabilistic encoder and Gaussian prior; the right panel shows the DAE pathway with swap noise corruption. Both pathways produce latent codes $\mathbf{z}$ and reconstruction errors $\boldsymbol{\epsilon}$ that are concatenated with the original features through the fusion hub for downstream stacking with gradient-boosted tree meta-learners.}
\label{fig:ae-architecture}
\end{figure}

\section{Experimental Environment Setup}

\subsection{Materials and Benchmark Dataset}

All experiments in this chapter are conducted on the Bank Account Fraud (BAF) Dataset Suite~\cite{feedzai2022baf,jesus2022bank}, published at NeurIPS 2022. The BAF dataset is a realistic, large-scale benchmark for tabular fraud detection derived from a real-world bank account opening process. It contains approximately one million instances across 30 raw features, with a fraud prevalence of 1.10\%---a class imbalance ratio of approximately 1:90. We specifically use Variant~I, which introduces controlled distributional shifts in feature space while maintaining the core data structure.

The BAF dataset is particularly well-suited for evaluating autoencoder-based approaches for several reasons. First, its large size ($\sim$1M instances) provides sufficient data to train deep autoencoder architectures without overfitting, while the temporal structure of the data (spanning 8 months) enables realistic train/validation/test splits that simulate production deployment conditions. Second, the moderate fraud rate of 1.10\% is representative of real-world fraud detection scenarios while being extreme enough to challenge supervised methods. Third, the heterogeneous feature set---comprising binary indicators, low-cardinality categoricals, and continuous numerical features---provides a realistic testbed for the tabular swap noise corruption mechanism. Fourth, the dataset has been carefully curated to remove directly identifying information, ensuring that models must learn behavioral patterns rather than memorizing identities.

The 30 raw features in the BAF dataset can be grouped into several semantic categories: demographic features (age, income, employment status), financial features (credit score, proposed credit limit, existing credit patterns), behavioral features (session length, device information, submission patterns), and temporal features (application month, days since last activity). After preprocessing and one-hot expansion of categorical features, the final feature dimension is $D = 53$. Table~\ref{tab:ae-feature-groups} summarizes the feature group composition after preprocessing.

\begin{table}[H]
\centering
\caption{Feature Group Summary After Preprocessing}
\label{tab:ae-feature-groups}
\small
\begin{tabularx}{0.95\textwidth}{lXrr}
\toprule
\textbf{Feature Group} & \textbf{Example Features} & \textbf{Raw Count} & \textbf{Expanded Dim} \\
\midrule
Demographic & income, date\_of\_birth, customer\_age & 5 & 9 \\
Financial & credit\_risk\_score, proposed\_credit\_limit & 7 & 12 \\
Behavioral & session\_length, device\_os, keep\_alive\_session & 8 & 18 \\
Temporal & month, days\_since\_last\_activity & 4 & 5 \\
Derived & velocity features, application count ratios & 6 & 9 \\
\midrule
\textbf{Total} & & \textbf{30} & \textbf{53} \\
\bottomrule
\end{tabularx}
\end{table}

\subsection{Preprocessing Pipeline}

The preprocessing pipeline transforms the raw BAF data into a clean, normalized feature matrix suitable for autoencoder training. The pipeline executes in six sequential phases, each designed to address a specific data quality or transformation requirement. Table~\ref{tab:ae-preprocessing-phases} provides a summary, and each phase is described in detail below.

\begin{table}[H]
\centering
\caption{Preprocessing Pipeline Phase Summary}
\label{tab:ae-preprocessing-phases}
\small
\begin{tabularx}{\textwidth}{clXl}
\toprule
\textbf{Phase} & \textbf{Name} & \textbf{Operation} & \textbf{Output} \\
\midrule
1 & Data Sanitation & Column name normalisation, type casting & Clean DataFrame \\
2 & Missing Value Imputation & Median (numeric), ``unknown'' (categorical) & Complete DataFrame \\
3 & Outlier Treatment & Winsorisation at 1st/99th percentiles & Clipped DataFrame \\
4 & Encoding and Scaling & Label encoding + MinMax normalisation & Numeric matrix \\
5 & Chronological Partitioning & Train: months 0--4, Val: month 5, Test: months 6--7 & Split tensors \\
6 & Vectorisation & DataFrame to PyTorch tensors & Training-ready data \\
\bottomrule
\end{tabularx}
\end{table}

\textbf{Phase~1---Data Sanitation:} Column names are sanitized by replacing non-alphanumeric characters with underscores and converting to lowercase. The target column (\texttt{fraud\_bool}) and temporal index (\texttt{month}) are identified and excluded from feature transformation to prevent label leakage and preserve the temporal partitioning structure. Feature types are explicitly validated and cast to their correct dtypes (integer for binary features, string for categorical features, float for continuous features).

\textbf{Phase~2---Missing Value Imputation:} Missing values in numerical columns are filled with the median of the corresponding training-set column, providing robustness to outliers that would bias the mean. Categorical missing values are assigned a dedicated ``unknown'' category, preserving the information that the value was missing---which can itself be a fraud indicator. Critically, imputation statistics are computed exclusively on the training set and applied uniformly to validation and test sets to prevent data leakage across temporal boundaries.

\textbf{Phase~3---Outlier Treatment:} Numerical features are winsorized at the 1st and 99th percentiles using training-set quantile boundaries. This clipping strategy mitigates the influence of extreme values that could destabilize the MinMax normalization and cause gradient instability during autoencoder training. The winsorization bounds are stored from the training set and applied consistently to validation and test sets.

\textbf{Phase~4---Encoding and Scaling:} Categorical features are one-hot encoded using training-set-fitted encoders, expanding the feature dimension from 30 to 53. Numerical features undergo MinMax normalization to the range $[0, 1]$, computed on the training set and clipped on validation and test sets. This produces a fully numerical feature matrix with all values in $[0, 1]$, suitable for both the sigmoid output activation of the autoencoders and the swap noise corruption of the DAE.

\textbf{Phase~5---Chronological Partitioning:} The dataset is partitioned temporally using the \texttt{month} column: months 0--4 constitute the training set, month 5 the validation set, and months 6--7 the test set. This temporal split simulates real-world deployment conditions where the model must generalize to future data distributions, and it avoids the optimistic bias that random splitting would introduce by allowing future information to leak into training. Importantly, for the target manifold isolation strategy, the autoencoder training set is further filtered to exclude all fraudulent transactions, ensuring that the model learns exclusively from the legitimate class.

\textbf{Phase~6---Vectorisation:} The preprocessed DataFrames are converted to PyTorch tensors with features stored as 32-bit floating-point arrays and labels as 64-bit integers. The resulting tensors are wrapped in PyTorch \texttt{TensorDataset} and \texttt{DataLoader} objects with configurable batch sizes for efficient GPU-based training. For the autoencoder training loop, the fraud-filtered legitimate-only training set is loaded without labels, while the full test set (including both classes) is used for evaluation.

\subsection{Hyperparameter Optimization}

The hyperparameter configurations for both the VAE and DAE models were determined through a systematic grid search guided by validation reconstruction loss. Table~\ref{tab:ae-hyperparameters} presents the final optimized configurations for both models, along with the rationale for each choice.

\begin{table}[H]
\centering
\caption{Hyperparameter Optimization Matrix: VAE vs.\ DAE}
\label{tab:ae-hyperparameters}
\small
\begin{tabularx}{\textwidth}{lXXX}
\toprule
\textbf{Hyperparameter} & \textbf{VAE} & \textbf{DAE} & \textbf{Rationale} \\
\midrule
Architecture & $53 \to 128 \to 64 \to 32 \to 64 \to 128 \to 53$ & $53 \to 128 \to 64 \to 32 \to 64 \to 128 \to 53$ & Symmetric hourglass; 32-dim bottleneck \\
Latent dimension ($d$) & 32 & 32 & Sufficient for manifold encoding; prevents over-compression \\
Encoder activations & ReLU & ReLU & Standard for tabular autoencoders \\
Decoder activations & ReLU + Sigmoid (output) & ReLU + Sigmoid (output) & Sigmoid maps output to $[0, 1]$ \\
Optimizer & Adam & Adam & Adaptive rates for sparse gradient regimes \\
Learning rate & $1 \times 10^{-3}$ & $1 \times 10^{-3}$ & Standard for Adam with reconstruction loss \\
Batch size & 512 & 512 & Balanced GPU utilisation and gradient stability \\
Epochs & 100 & 100 & Sufficient for convergence with early stopping \\
Early stopping patience & 10 & 10 & Based on validation reconstruction loss \\
KL weight ($\beta$) & 1.0 & N/A & Standard $\beta$-VAE with $\beta=1$ \\
Swap noise rate ($\nu$) & N/A & 0.15 & Marginal-preserving corruption; 15\% of features \\
Loss function & ELBO (MSE + KL) & MSE & VAE: probabilistic; DAE: deterministic \\
Training data & Legitimate only & Legitimate only & Target manifold isolation \\
Weight initialisation & Xavier & Xavier & Stable gradient flow in symmetric architecture \\
Dropout & 0.0 & 0.0 & Swap noise provides regularisation for DAE \\
\bottomrule
\end{tabularx}
\end{table}

Several design choices warrant further discussion. The latent dimension of 32 was selected based on a preliminary search over $\{8, 16, 32, 64\}$, which showed that $d = 32$ provides the best trade-off between reconstruction fidelity and manifold compression. Smaller latent dimensions ($d \leq 16$) produced poor reconstructions even for legitimate transactions, while larger dimensions ($d = 64$) increased the risk of the identity mapping problem by providing the autoencoder with sufficient capacity to reconstruct arbitrary inputs. The corruption rate $\nu = 0.15$ for the DAE was selected from a search over $\{0.05, 0.10, 0.15, 0.20, 0.30\}$, with $\nu = 0.15$ yielding the best validation anomaly separation (highest AUC for standalone DAE scoring). Lower corruption rates provided insufficient regularisation, while higher rates degraded reconstruction quality to the point where the anomaly signal was overwhelmed by reconstruction noise.

The KL weight $\beta = 1.0$ follows the standard $\beta$-VAE formulation, providing equal weighting between the reconstruction and regularisation terms. We experimented with $\beta \in \{0.1, 0.5, 1.0, 2.0, 5.0\}$ and found that lower $\beta$ values (stronger reconstruction emphasis) led to near-deterministic latent codes that overfit the training distribution, while higher $\beta$ values (stronger regularisation) caused posterior collapse where the encoder ignored the input. The standard $\beta = 1.0$ provided the best balance between manifold regularisation and reconstruction fidelity.

\begin{figure}[H]
\centering
\includegraphics[width=0.85\textwidth]{ae_fig_01}
\caption{Hyperparameter sensitivity analysis. Left: Effect of latent dimension $d$ on reconstruction loss and anomaly separation AUC for both VAE and DAE. Right: Effect of swap noise rate $\nu$ on DAE validation performance, showing the trade-off between regularisation strength and reconstruction quality.}
\label{fig:ae-hyperparameter-sensitivity}
\end{figure}

\section{Experimental Results}

\subsection{The Deceptive Properties of Global ROC Metrics}

Before presenting the detailed experimental results, it is essential to address a methodological subtlety that has significant implications for interpreting anomaly detection performance under extreme class imbalance: the deceptive nature of global ROC-AUC as an evaluation metric. The ROC curve plots the true positive rate (TPR) against the false positive rate (FPR) as the classification threshold is varied, and the area under this curve (AUC) summarizes the model's ranking ability across all possible operating thresholds. While AUC is threshold-independent and widely used, its interpretation under extreme class imbalance can be profoundly misleading.

Consider a fraud detection scenario with 1.10\% fraud prevalence. At an FPR of 5\%, which is already a generous operational budget for a production system, the model has access to only the leftmost 5\% of the FPR axis on the ROC curve. The remaining 95\% of the curve corresponds to operating points with FPR $> 5\%$, which are operationally irrelevant because they would generate an unacceptable volume of false alerts. Yet these irrelevant operating points contribute 95\% of the area under the ROC curve. Consequently, a model can achieve a deceptively high AUC by performing well on the irrelevant high-FPR portion of the curve while performing poorly on the operationally critical low-FPR region.

This phenomenon is particularly acute for autoencoder-based anomaly detectors. The standalone VAE achieves an AUC of 0.5501, which appears to be only marginally above the random baseline of 0.5. However, this near-random AUC obscures the fact that the VAE does capture some meaningful anomaly signal---it simply fails to separate the classes across the full range of thresholds. When the VAE's anomaly features are combined with supervised meta-learners through stacking, the resulting model achieves an AUC of 0.8868, demonstrating that the anomaly signal was present but too weak and noisy for threshold-based discrimination. The global AUC metric, by averaging over all thresholds including operationally irrelevant ones, fails to capture this nuance and can lead to the premature dismissal of autoencoder-derived features that provide valuable complementary information in a stacking framework.

We therefore supplement the global AUC with operationally grounded metrics: TPR at fixed FPR thresholds (5\%, 10\%, 20\%), precision at these thresholds, and the partial AUC computed over the $[0, 0.05]$ FPR interval. These metrics provide a more faithful picture of model performance under realistic deployment constraints.

\subsection{Systematic Performance Consolidation}

Table~\ref{tab:ae-main-results} presents the systematic performance consolidation across all model configurations evaluated in this chapter. The results are organized into three tiers: standalone unsupervised models (VAE, DAE), standalone supervised models (LightGBM, XGBoost), and hybrid stacked models that combine unsupervised autoencoder features with supervised meta-learners.

\begin{table}[H]
\centering
\caption{Systematic Performance Consolidation: All Model Configurations}
\label{tab:ae-main-results}
\small
\setlength{\tabcolsep}{5pt}
\begin{tabularx}{\textwidth}{lXXXXX}
\toprule
\textbf{Model} & \textbf{AUC} & \textbf{TPR @ 5\% FPR} & \textbf{Precision @ 5\% FPR} & \textbf{F1 (Fraud)} & \textbf{Type} \\
\midrule
VAE (standalone) & 0.5501 & 9.94\% & 0.024 & 0.038 & Unsupervised \\
DAE (standalone) & 0.5210 & 6.30\% & 0.023 & 0.032 & Unsupervised \\
\midrule
LightGBM (baseline) & 0.8798 & 48.52\% & 0.116 & 0.187 & Supervised \\
XGBoost (baseline) & 0.8830 & 50.13\% & 0.119 & 0.193 & Supervised \\
\midrule
DAE + XGBoost (stacked) & 0.8752 & 49.80\% & 0.118 & 0.191 & Hybrid \\
VAE + Tree Ensemble (stacked) & 0.8868 & 52.99\% & 0.124 & 0.201 & Hybrid \\
\bottomrule
\end{tabularx}
\end{table}

The results reveal a striking performance dichotomy between standalone and stacked configurations. Standalone autoencoder anomaly scores achieve AUC values only marginally above random guessing (0.5501 for VAE, 0.5210 for DAE), with TPR values at 5\% FPR of just 9.94\% and 6.30\% respectively. These results confirm the well-known limitation of reconstruction-based anomaly detection: the autoencoder reconstructs both legitimate and fraudulent transactions with similar fidelity, providing insufficient discrimination for threshold-based anomaly scoring. The DAE performs even worse than the VAE in the standalone setting, suggesting that the deterministic reconstruction with swap noise regularisation produces slightly better reconstruction of out-of-distribution inputs compared to the probabilistic VAE decoder, further collapsing the anomaly signal.

However, the stacked configurations tell a dramatically different story. The VAE + Tree Ensemble model achieves an AUC of 0.8868 and a TPR of 52.99\% at 5\% FPR, surpassing both the standalone LightGBM baseline (AUC 0.8798, TPR 48.52\%) and the standalone XGBoost baseline (AUC 0.8830, TPR 50.13\%). The DAE + XGBoost stacked model achieves an AUC of 0.8752 and a TPR of 49.80\%, which is competitive with but slightly below the supervised baselines. The VAE + Tree Ensemble configuration represents the best-performing model in this chapter, demonstrating that the probabilistic latent regularisation of the VAE provides more useful features for downstream stacking than the deterministic DAE reconstruction.

\begin{figure}[H]
\centering
\includegraphics[width=0.92\textwidth]{ae_fig_02}
\caption{Systematic performance comparison across all model configurations. Left: ROC-AUC for standalone, supervised, and stacked models. Center: TPR at 5\% FPR, the operationally critical metric. Right: Fraud F1-score, reflecting the harmonic mean of precision and recall for the minority class. The VAE + Tree Ensemble configuration achieves the highest performance across all three metrics.}
\label{fig:ae-performance-comparison}
\end{figure}

\subsection{VAE Continuous Latent Embedding Performance}

The Variational Autoencoder produces two types of anomaly features: the 32-dimensional latent code $\mathbf{z}_i$ and the 53-dimensional element-wise reconstruction error $\boldsymbol{\epsilon}_i$. In the standalone setting, the aggregate reconstruction error $\|\boldsymbol{\epsilon}_i\|$ is used as the anomaly score. Table~\ref{tab:ae-vae-results} presents the detailed performance metrics for the VAE in both standalone and stacked configurations.

\begin{table}[H]
\centering
\caption{VAE Performance: Standalone vs.\ Stacked}
\label{tab:ae-vae-results}
\small
\begin{tabularx}{0.95\textwidth}{lXXXXXXXX}
\toprule
\textbf{Config.} & \textbf{AUC} & \textbf{TPR @5\%} & \textbf{TPR @10\%} & \textbf{Prec. @5\%} & \textbf{F1} & \textbf{TN} & \textbf{FP} & \textbf{TP} \\
\midrule
Standalone VAE & 0.5501 & 9.94\% & 15.67\% & 0.024 & 0.038 & 192,013 & 10,120 & 286 \\
VAE + Tree Ens. & 0.8868 & 52.99\% & 67.12\% & 0.124 & 0.201 & 192,008 & 10,125 & 1,525 \\
\bottomrule
\end{tabularx}
\end{table}

The standalone VAE achieves an AUC of only 0.5501, barely above the random baseline of 0.5. At the 5\% FPR operating point, the TPR of 9.94\% means that the VAE identifies fewer than 1 in 10 fraudulent transactions while generating over 10,000 false alerts. This poor performance is attributable to several factors. First, the VAE learns to reconstruct the dominant mode of the data distribution (legitimate transactions) with high fidelity, but the probabilistic latent structure does not create a sufficiently sharp boundary between in-distribution and out-of-distribution inputs. The Gaussian prior encourages smooth interpolation across the latent space, which paradoxically enables the decoder to produce reasonable reconstructions for inputs that lie between the dense regions of the legitimate manifold, including many fraudulent transactions that fall in the tails of the legitimate distribution. Second, the ELBO objective balances reconstruction accuracy against KL regularisation, and this balance point may not align with the optimal point for anomaly detection. The KL term pushes the encoder toward the prior, which is centred at the origin and may assign overlapping latent codes to legitimate and fraudulent inputs that share similar feature values.

The dramatic improvement to AUC 0.8868 in the stacked configuration (a $\Delta$AUC of +0.3367) reveals that the VAE's latent features and reconstruction errors, while insufficient for standalone anomaly detection, provide highly complementary information when combined with supervised classification. The 32-dimensional latent code captures a compressed representation of each transaction's position on the learned manifold, and the gradient-boosted trees can learn decision boundaries in this latent space that are not accessible through the original features alone. The reconstruction error profile provides a feature-level anomaly signal that, when weighted and combined by the meta-learner, contributes to improved discrimination.

\begin{figure}[H]
\centering
\includegraphics[width=0.85\textwidth]{ae_fig_03}
\caption{VAE latent space visualization. Left: 2D t-SNE projection of the 32-dimensional VAE latent codes, with legitimate transactions shown in blue and fraudulent transactions in red. The fraud instances are not cleanly separated from the legitimate cluster, explaining the poor standalone performance. Right: Distribution of aggregate reconstruction error $\|\boldsymbol{\epsilon}\|$ for legitimate and fraudulent transactions, showing substantial overlap between the two distributions.}
\label{fig:ae-vae-latent}
\end{figure}

\begin{figure}[H]
\centering
\includegraphics[width=0.85\textwidth]{ae_fig_04}
\caption{VAE ROC curves. Left: Standalone VAE anomaly scoring (AUC = 0.5501). Right: VAE + Tree Ensemble stacked model (AUC = 0.8868). The shaded region indicates the operationally critical FPR $\leq$ 5\% zone. The dramatic AUC improvement in the stacked model demonstrates that VAE-derived features provide valuable complementary information for the supervised meta-learner.}
\label{fig:ae-vae-roc}
\end{figure}

\subsection{DAE Deterministic Reconstruction Performance}

The Denoising Autoencoder produces the same two types of anomaly features as the VAE---latent codes and reconstruction errors---but through a fundamentally different mechanism. The DAE's deterministic encoder maps each input to a single latent code, and the swap noise regularisation forces the network to learn robust features that capture the joint distribution of the training data. Table~\ref{tab:ae-dae-results} presents the detailed performance metrics.

\begin{table}[H]
\centering
\caption{DAE Performance: Standalone vs.\ Stacked}
\label{tab:ae-dae-results}
\small
\begin{tabularx}{0.95\textwidth}{lXXXXXXXX}
\toprule
\textbf{Config.} & \textbf{AUC} & \textbf{TPR @5\%} & \textbf{TPR @10\%} & \textbf{Prec. @5\%} & \textbf{F1} & \textbf{TN} & \textbf{FP} & \textbf{TP} \\
\midrule
Standalone DAE & 0.5210 & 6.30\% & 10.89\% & 0.023 & 0.032 & 192,019 & 10,114 & 181 \\
DAE + XGBoost & 0.8752 & 49.80\% & 63.45\% & 0.118 & 0.191 & 192,011 & 10,122 & 1,434 \\
\bottomrule
\end{tabularx}
\end{table}

The standalone DAE achieves an AUC of 0.5210, even lower than the standalone VAE. At the 5\% FPR operating point, the TPR of 6.30\% indicates that the DAE identifies barely 1 in 15 fraudulent transactions. This worse performance relative to the VAE can be attributed to the deterministic nature of the DAE's latent representation. Without the stochastic sampling of the VAE, the DAE's encoder maps each input to a single point in latent space, and the decoder must reconstruct the input from this point alone. For legitimate inputs, this is straightforward because the DAE has been trained on this distribution. However, for fraudulent inputs that share many feature values with legitimate ones (as is typical in fraud detection), the deterministic encoder maps them to nearby latent codes, and the decoder produces reconstructions that are similar to what it would produce for the corresponding legitimate inputs. The result is low reconstruction error for fraudulent transactions that are ``close'' to the legitimate manifold, collapsing the anomaly signal.

The swap noise regularisation does provide some benefit: by training the DAE to recover original values from corrupted inputs, the network learns cross-feature dependencies that enable it to detect transactions with unusual feature combinations. However, this benefit is limited because the swap noise corruption preserves marginal feature statistics---the corrupted values are always legitimate values drawn from the training distribution. Consequently, the DAE's reconstruction quality for corrupted inputs is high, and the reconstruction error for fraudulent inputs is only marginally elevated relative to the error for legitimate inputs.

The stacked DAE + XGBoost configuration achieves an AUC of 0.8752, which is below both the standalone supervised baselines and the stacked VAE configuration. The $\Delta$AUC improvement from stacking (+0.3542) is actually larger than for the VAE (+0.3367), indicating that the DAE's features provide substantial complementary information. However, the starting baseline is lower (0.5210 vs.\ 0.5501), and the absolute stacked performance falls short. The DAE's deterministic reconstruction errors, while informative about feature-level anomalies, appear to provide a noisier signal than the VAE's probabilistic reconstruction errors, which incorporate the uncertainty encoded in the latent distribution.

\begin{figure}[H]
\centering

\includegraphics[width=0.85\textwidth]{ae_fig_05}
\caption{DAE reconstruction error analysis. Left: Per-feature average reconstruction error for legitimate (blue) and fraudulent (red) transactions. The error profiles are highly correlated across features, with fraud showing only marginal elevation. Right: Scatter plot of aggregate reconstruction error vs.\ latent code norm, showing the overlap between legitimate and fraudulent transactions in this two-dimensional anomaly space.}
\label{fig:ae-dae-errors}
\end{figure}

\begin{figure}[H]
\centering
\includegraphics[width=0.85\textwidth]{ae_fig_06}
\caption{DAE ROC curves. Left: Standalone DAE anomaly scoring (AUC = 0.5210). Right: DAE + XGBoost stacked model (AUC = 0.8752). The stacked DAE model performs competitively with supervised baselines but does not surpass them, unlike the VAE + Tree Ensemble configuration.}
\label{fig:ae-dae-roc}
\end{figure}

\subsection{The Resampling and Miscalibration Paradox}

A natural question arising from the target manifold isolation strategy is whether resampling techniques such as SMOTE~\cite{chawla2002smote} or ADASYN~\cite{he2008adasyn} could improve autoencoder performance by augmenting the fraud class during training. To investigate this, we conducted additional experiments where the autoencoder training set included fraudulent transactions alongside synthetic fraud samples generated by SMOTE and ADASYN. The results were uniformly negative: including fraud samples in the autoencoder training set degraded standalone anomaly detection performance and, more critically, introduced systematic miscalibration in the stacked models.

The mechanism behind this miscalibration paradox is straightforward. SMOTE and ADASYN generate synthetic minority samples by interpolating between existing fraud instances in feature space. These synthetic samples lie on the linear segments connecting real fraud instances, which are themselves embedded within the cloud of legitimate transactions. When the autoencoder is trained on these synthetic fraud samples, it learns to reconstruct them accurately, effectively extending the ``normal manifold'' to encompass the interpolated fraud regions. This manifold extension has two deleterious effects. First, it reduces the reconstruction error for real fraud instances that lie near the interpolated regions, further collapsing the anomaly signal. Second, it creates a distorted reconstruction error landscape where the error is elevated not only for genuine anomalies but also for legitimate transactions that lie near the synthetic sample boundaries, increasing the false positive rate.

Table~\ref{tab:ae-resampling-paradox} summarizes the performance degradation observed when resampling is applied inside the autoencoder training loop.

\begin{table}[H]
\centering
\caption{The Resampling Miscalibration Paradox: Performance Degradation with SMOTE/ADASYN}
\label{tab:ae-resampling-paradox}
\small
\begin{tabularx}{0.95\textwidth}{lXXXX}
\toprule
\textbf{Configuration} & \textbf{Standalone AUC} & \textbf{Stacked AUC} & \textbf{TPR @ 5\% FPR} & \textbf{Calibration Error} \\
\midrule
VAE (legitimate only) & 0.5501 & 0.8868 & 52.99\% & 0.031 \\
VAE + SMOTE & 0.5312 & 0.8719 & 47.64\% & 0.089 \\
VAE + ADASYN & 0.5278 & 0.8694 & 46.83\% & 0.097 \\
\midrule
DAE (legitimate only) & 0.5210 & 0.8752 & 49.80\% & 0.035 \\
DAE + SMOTE & 0.5089 & 0.8614 & 44.51\% & 0.102 \\
DAE + ADASYN & 0.5063 & 0.8587 & 43.29\% & 0.114 \\
\bottomrule
\end{tabularx}
\end{table}

The calibration error, measured as the expected absolute difference between predicted probabilities and observed frequencies, increases by a factor of 2.5--3.3$\times$ when resampling is applied. This miscalibration is particularly problematic in production fraud detection systems, where the predicted probability must be interpretable as the likelihood of fraud: a miscalibrated model that assigns systematically higher or lower probabilities than warranted by the evidence can lead to either excessive manual review costs (over-alerting) or missed fraud (under-alerting). The target manifold isolation strategy---training exclusively on legitimate transactions---avoids this miscalibration by maintaining a clean separation between the unsupervised feature extraction stage (autoencoder) and the supervised classification stage (meta-learner), ensuring that each component operates within its theoretical competence.

\begin{figure}[H]
\centering
\includegraphics[width=0.92\textwidth]{ae_fig_07}
\caption{The resampling miscalibration paradox. Left: Reliability diagrams comparing calibration of VAE + Tree Ensemble models trained with legitimate-only data versus SMOTE-augmented data. The SMOTE model shows systematic over-prediction at lower probability bins. Right: TPR at 5\% FPR across resampling configurations, demonstrating consistent degradation when synthetic samples are included in autoencoder training.}
\label{fig:ae-resampling-paradox}
\end{figure}

\begin{figure}[H]
\centering
\includegraphics[width=0.85\textwidth]{ae_fig_08}
\caption{Stacking benefit decomposition. The bar chart shows the AUC contribution of each feature source in the VAE + Tree Ensemble model: original features alone (blue), latent codes alone (green), reconstruction errors alone (red), and the full concatenated feature set (purple). The full stacking configuration outperforms any individual feature source, confirming that the autoencoder-derived features provide genuinely complementary information.}
\label{fig:ae-stacking-decomposition}
\end{figure}

\begin{figure}[H]
\centering
\includegraphics[width=0.85\textwidth]{ae_fig_09}
\caption{Feature importance in the VAE + Tree Ensemble stacked model. The top 20 features by XGBoost gain importance include both original features and autoencoder-derived features (reconstruction errors and latent dimensions), confirming that the meta-learner actively utilizes the unsupervised anomaly signal.}
\label{fig:ae-feature-importance}
\end{figure}

\begin{figure}[H]
\centering

\includegraphics[width=0.85\textwidth]{ae_fig_10}
\caption{Per-feature reconstruction error comparison between VAE and DAE. The heatmap shows the average element-wise reconstruction error $\epsilon_j = |x_j - \hat{x}_j|$ for each of the 53 features, separately for legitimate and fraudulent transactions. Both models show elevated errors for fraud, but the signal-to-noise ratio is insufficient for standalone discrimination.}
\label{fig:ae-per-feature-errors}
\end{figure}

\section{Chapter Summary}

This chapter has presented Pipeline~III, an unsupervised feature extraction framework based on Variational and Denoising Autoencoders for financial fraud detection. The pipeline addresses the inductive bias failure of supervised classifiers under extreme class imbalance by learning a reconstruction-based model of the legitimate transaction manifold and using deviations from this manifold as anomaly features for downstream classification.

The key findings are as follows:

\begin{enumerate}
\item \textbf{Standalone autoencoder anomaly scoring is insufficient.} Both the VAE (AUC 0.5501) and DAE (AUC 0.5210) produce anomaly scores that barely exceed random guessing when used in isolation. The fundamental limitation is that autoencoders trained on legitimate transactions can still reconstruct fraudulent transactions with relatively low error, particularly when the fraud instances lie in the tails of the legitimate distribution rather than in a clearly separated region of feature space.

\item \textbf{Score-level stacking provides dramatic improvement.} When autoencoder-derived features (latent codes $\mathbf{z}$ and element-wise reconstruction errors $\boldsymbol{\epsilon}$) are concatenated with the original features and fed to gradient-boosted tree meta-learners, performance improves substantially. The VAE + Tree Ensemble achieves AUC 0.8868 (TPR 52.99\% at 5\% FPR), surpassing both standalone supervised baselines (LightGBM AUC 0.8798, XGBoost AUC 0.8830). The DAE + XGBoost achieves AUC 0.8752 (TPR 49.80\%), which is competitive but does not exceed the supervised baselines.

\item \textbf{VAE features outperform DAE features for stacking.} The probabilistic latent structure of the VAE provides more useful features for downstream stacking than the deterministic DAE reconstruction. The VAE's stochastic sampling introduces uncertainty into the latent codes that appears to regularize the meta-learner's feature utilization, while the DAE's deterministic mapping produces reconstruction errors that, while informative, provide a noisier anomaly signal.

\item \textbf{Target manifold isolation is critical.} Training the autoencoders exclusively on legitimate transactions is essential for creating a meaningful anomaly signal. Including fraudulent or synthetic samples in the autoencoder training set extends the learned manifold to encompass fraud patterns, reducing reconstruction errors for fraudulent transactions and collapsing the anomaly discrimination.

\item \textbf{Resampling inside autoencoder loops causes miscalibration.} Applying SMOTE or ADASYN to augment the fraud class during autoencoder training degrades both standalone and stacked performance while introducing systematic probability miscalibration. The calibration error increases by 2.5--3.3$\times$ compared to the legitimate-only training strategy, making the resulting models unreliable for production deployment.
\end{enumerate}

These findings demonstrate that autoencoder-based unsupervised feature extraction, while insufficient as a standalone fraud detection mechanism, provides valuable complementary information that enhances supervised classifiers when integrated through a principled stacking framework. The VAE + Tree Ensemble configuration represents the best-performing model in this pipeline, achieving state-of-the-art performance on the BAF Dataset Variant~I by combining the distributional modelling capabilities of the VAE with the discriminative power of gradient-boosted decision trees. The results also highlight the importance of careful experimental design---particularly the avoidance of resampling within the autoencoder training loop---and the need for operationally grounded evaluation metrics that look beyond global AUC to assess performance under realistic deployment constraints.
